% Part: propositional-logic
% Chapter: syntax-and-semantics
% Section: introduction

\documentclass[../../../include/open-logic-section]{subfiles}

\begin{document}

\olfileid[fr]{pl}{syn}{int}

\olsection{Introduction}

La logique propositionnelle étudie des !!{formula}s construites à partir
de !!{propositional variable}s au moyen des connecteurs propositionnels
$\lnot$, $\land$, $\lor$, $\lif$ et $\liff$. Intuitivement,
!!a{propositional variable}~$p$ représente une phrase ou une proposition
qui est vraie ou fausse. Dès que les valeurs de vérité des
!!{propositional variable}s sont fixées, celles de toutes les !!{formula}s
construites à partir d'elles au moyen de ces connecteurs le sont aussi.
On dit que la logique propositionnelle est \emph{vérifonctionnelle},
car sa sémantique est donnée par des fonctions des valeurs de vérité.
En particulier, on laisse de côté toute autre détermination du vrai
et du faux : on ne se demande pas si quelque chose est nécessairement
vrai plutôt que vrai de façon contingente, si l'on sait que c'est vrai,
ou si c'est vrai maintenant plutôt que dans le passé ou dans l'avenir.
On ne considère que deux valeurs de vérité, le vrai ($\True$) et le faux
($\False$), et l'on exclut donc de cette étude la possibilité qu'un
énoncé ne soit ni vrai ni faux, ou ne soit qu'à moitié vrai. On se
limite aussi aux connecteurs pour lesquels la valeur de vérité d'une
!!{formula} composée dépend entièrement des valeurs de vérité de ses
constituants, et non, par exemple, de leur sens. En particulier, le
caractère vérifonctionnel des conditionnels de l'anglais, en ce sens,
est controversé. L'implication matérielle~$\lif$ est vérifonctionnelle ;
d'autres logiques étudient des conditionnels qui ne le sont pas.

Pour développer la théorie et la métathéorie de la logique
propositionnelle vérifonctionnelle, il faut d'abord définir la syntaxe
et la sémantique de ses expressions. Nous présenterons une manière de
construire des !!{formula}s à partir de !!{propositional variable}s au
moyen des connecteurs. D'autres définitions sont possibles. D'autres
systèmes emploient d'autres symboles, prennent d'autres ensembles de
connecteurs comme primitifs et utilisent autrement les parenthèses,
voire s'en passent, comme dans la notation dite polonaise. Toutes ces
approches ont cependant en commun de définir l'ensemble des !!{formula}s
par des règles de formation \emph{inductives}. Si ces règles sont bien
conçues, chaque expression ne peut essentiellement être construite que
d'une seule manière. Cette définition inductive d'expressions à
\emph{lecture unique} permet de leur attribuer un sens par le même
procédé : une définition inductive.

L'attribution d'un sens aux expressions relève de la sémantique. En
logique propositionnelle, la notion centrale est celle de satisfaction
par une !!{valuation}. Une !!{valuation}~$\pAssign{v}$ attribue les
valeurs de vérité $\True$, $\False$ aux !!{propositional variable}s.
Toute !!{valuation} détermine une valeur de vérité $\pValue{v}(!A)$
pour chaque !!{formula}~$!A$. Une !!{formula} est satisfaite par une
!!{valuation}~$\pAssign{v}$ si et seulement si
$\pValue{v}(!A) = \True$ ; on écrit alors $\pSat{v}{!A}$.
Cette relation peut elle aussi se définir par induction sur la
structure de~$!A$. Les fonctions de vérité des connecteurs permettent
par exemple de définir la satisfaction de $!A \land !B$ à partir de
la satisfaction, ou de la non-satisfaction, de $!A$ et de~$!B$.

À partir de la relation de satisfaction $\pSat{v}{!A}$ pour les
formules, on définit les notions sémantiques fondamentales de
tautologie, de conséquence logique et de satisfaisabilité.
Une !!{formula} est une tautologie, $\Entails !A$, si toute
!!{valuation} la satisfait, c'est-à-dire si $\pValue{v}(!A) = \True$
pour toute~$\pAssign{v}$. Elle est conséquence logique d'un ensemble
de !!{formula}s, $\Gamma \Entails !A$, si toute !!{valuation} qui
satisfait toutes les !!{formula}s de~$\Gamma$ satisfait également~$!A$.
Un ensemble de !!{formula}s est satisfaisable s'il existe une
!!{valuation} qui satisfait simultanément toutes ses !!{formula}s.
Les !!{formula}s étant définies inductivement, et la satisfaction
étant à son tour définie par induction sur leur structure, on peut
utiliser l'induction pour démontrer des propriétés de cette sémantique
et établir des liens entre les notions ainsi définies.

\end{document}
